\section{Sammendrag}
Vi har satt oss fore å lage en maskin som viser vår kunnskap om mekanismer, PCB-design og programmering. Hovedsystemet er satt sammen av en Raspberry Pi, tilpassede PCB-er, 3D-printede deler og andre elektronikkomponenter. Vi har tilpasset alle komponenter slik at de passer sammen, og designet PCB-ene for å levere den nødvendige funksjonaliteten: PWM-utganger, I/O-knapper, HX711 loadcell-brikker og dedikerte I2C-tilkoblinger. Programvaren gir PID-regulering, IK-posisjonering og tilstandsmaskiner. Til slutt har vi bygget, testet, feilrettet og ferdigstilt en robotarm som fyller glass når vektsensoren registrerer at de er tomme. Vi har dermed nådd målet for semesterprosjektet.